\section{系统运行与复现说明}

\subsection{运行环境}
\begin{itemize}
\item Python 3.11(conda 环境),Go 1.26(终端客户端);PostgreSQL 16 + pgvector 扩展。
\item 关键依赖:\texttt{fastapi}、\texttt{psycopg}、\texttt{pgvector}、\texttt{sentence-transformers}、\texttt{scikit-learn}、\texttt{networkx}、\texttt{pandas}。
\item 本地大模型(可选):vLLM-Metal / LM Studio 等 OpenAI 兼容服务;无则自动用 mock,检索证据仍真实。
\end{itemize}

\subsection{一键安装}
\begin{codepanel}\ttfamily\small
make install\hfill\# 安装后端 Python 依赖
\end{codepanel}

\subsection{构建服务层与模型文件}
\begin{codepanel}\ttfamily\small
\# 1. 服务层(建库 + 装载语料/片段)\\
createdb sciscope\\
make postgres-schema POSTGRES\_DSN=postgresql://USER@localhost:5432/sciscope\\
psql sciscope -f infra/postgres/pgvector.sql\\
make postgres-load POSTGRES\_DSN=...\\[3pt]
\# 2. 数据治理:去重(先 dry-run 再 APPLY)\\
make dedupe-db\hfill\# 干跑统计\\
make dedupe-db APPLY=1\hfill\# 实际去重\\[3pt]
\# 3. 模型层:嵌入 + 推荐 + 趋势 + 图谱\\
SCISCOPE\_EMBEDDER\_PATH=models/embedder\_local/multilingual-e5-base make embeddings\\
make recommend-model trend-model graph-export\\
\# 或一键:make agent-build
\end{codepanel}

\subsection{启动与体验}
\begin{codepanel}\ttfamily\small
\# 带数据库 + 本地嵌入启动后端\\
SCISCOPE\_DB\_DSN=postgresql://USER@localhost:5432/sciscope \textbackslash\\
SCISCOPE\_EMBEDDER\_PATH=models/embedder\_local/multilingual-e5-base make dev\\[3pt]
\# 终端科研智能体(Go 客户端;本地开发需后端 :8000 与本地大模型 :8001,\\
\# 演示服务器固定 :8010,SCISCOPE\_BACKEND=http://10.21.16.192:8010 直连)\\
make backend \&  make llm \&  make tui\\[3pt]
\# 接本地大模型生成(需 :8001 上的 OpenAI 兼容服务)\\
make vllm-serve\hfill\# 启动本地模型\\
... make dev\hfill\# 配 LOCAL\_LLM\_* 指向该模型
\end{codepanel}
使用 \texttt{make tui} 进入终端科研智能体;后端接口文档:本地 \path{http://127.0.0.1:8000/docs},演示服务器 \path{http://10.21.16.192:8010/docs}。

\subsection{预编译客户端安装}
若只体验终端客户端,可直接安装预编译的 Go TUI,再连接公开 demo 后端或自建后端:
\begin{codepanel}\ttfamily\small
\# 跨平台开发者路径(Node.js)\\
npm install -g sciscope-tui\\[3pt]
\# macOS(Homebrew)\\
brew tap Timcai06/sciscope\\
brew trust --cask timcai06/sciscope/sciscope-tui\\
brew install --cask sciscope-tui\\[3pt]
\# Windows(Scoop)\\
scoop bucket add sciscope https://github.com/Timcai06/scoop-sciscope\\
scoop install sciscope-tui
\end{codepanel}
上述安装只分发终端客户端本体;若要体现全量科研智能体能力,仍需连接后端服务。公开后端适合演示和轻量试用,全量语料与长期运行建议按本章前述步骤部署本地或自建后端。

\subsection{质量与效果验证}
\begin{codepanel}\ttfamily\small
make test-backend\hfill\# 后端自动化测试(以当前环境输出为准)\\
make agent-smoke\hfill\# live agent 黑盒验收:库表规模、skills、工具预算\\
make eval-retrieval\hfill\# 自检索评估:recall@k / MRR / 延迟
\end{codepanel}

\begin{plainbox}
\textbf{复现边界(据实)}:本机开发基线与远端 2080 Ti 运行态需要分开看待。2026-08-10 已验证的远端全量运行态包括
\texttt{chunk\_embeddings=367{,}861/367{,}861}、\texttt{paper\_embeddings=159{,}164/159{,}164} 与 ANN 索引; 2026-08-14 演示服务器完成讯飞交付包 5{,}566 篇全链路入库(\texttt{papers=164{,}730}、\texttt{chunk\_embeddings=413{,}971}),后端 0.0.0.0:8010 对外服务。本机报告构建与接口调试可不携带这组全量资产。公开托管后端为 demo-lite 子集,不能替代本地/自建全量部署的能力证明。
\end{plainbox}

\subsection{全量重建(数据更新后)}
持续爬取积累新数据后,一次性重建并自动并入去重与全部模型:
\begin{codepanel}\ttfamily\small
make data-layer-refresh \&\& make rag-chunks \&\& make postgres-refresh \&\& make agent-build
\end{codepanel}

\subsection{交付物清单}
\begin{itemize}
\item \textbf{代码}:\path{src/}(采集/治理/分析/模型)、\path{backend/}(API + 智能体编排)、\path{tui/}(Go 终端客户端)、\path{infra/}(schema)、\path{Makefile}(流水线)。
\item \textbf{模型文件}:pgvector 向量索引、\texttt{models/}(嵌入器/趋势/推荐)、\texttt{output/graphs/}(图谱 JSON)、RAG 片段语料;均可由 \texttt{make agent-build} 复现。
\item \textbf{报告}:本《项目报告书》与同级《数据分析报告》(\texttt{output/pdf/sciscope\_data\_report})。
\item \textbf{演示样例}:\texttt{docs/examples/golden\_verify\_claim\_session.md} 给出一轮跨语言论断核查黄金会话,评委即使不启动环境也可查看工作流阶段、证据时间线与最终结论格式。
\item \textbf{黑盒验收}:\texttt{make agent-smoke} 调用 live \texttt{/api/agent},检查数据库级语料规模、能力边界、论断核查、趋势分析与论文推荐工具链路。
\item \textbf{基础大模型}:声明为可替换依赖(任意 OpenAI 兼容端点),不绑定特定模型。
\item \textbf{提交附件边界}:为满足逐文件上传与单文件大小限制,正式附件提供源代码、运行说明、评测/图谱/小型模型结果和校验清单;大型数据目录、嵌入器权重与数据库运行表不随附件上传,由 \texttt{Makefile} 与运行说明重建。
\end{itemize}

\vspace{4mm}
\begin{center}\small\color{SciMuted}
SciScope · 面向科技文献智能分析的科研智能体 · 大模型可换,证据可溯,流水线可复现
\end{center}
